(1) 封闭定容、隔热且完全混匀的反应器内，初始 \(N_0\) 摩尔物质 \(A\) 按 \(A\to B\) 发生一级转化，速率 \(\dot N_A=-kN_A\)，常数 \(k>0\)，每摩尔释放热量 \(Q>0\)，整体有效定容热容为常数 \(C>0\)。求浓度变化与转化率。

(2) 建立反应热使温度变化的能量账本并求 \(T(t)\)、最终绝热温度；解释温升为何有上界。

(3) 求瞬时热释放功率及其峰值和累计热量、温升到总温升一半的时刻。

(4) 如果实际速率为Arrhenius型 \(k(T)=k_0e^{-E_a/(RT)}\)，建立两条耦合ODE并证明能量约束 \(T+(Q/C)N_A\) 为常数；说明为何此时一般不能直接使用第一问指数解。
